% ==============================================================================
% Fichier     : main.tex
% Titre       : Supervision, Analyse Forensique et Gestion des Incidents
% Auteur      : MINKA MI NGUIDJOÏ Thierry Emmanuel
% Institution : Université de Yaoundé I — Faculté des Sciences
%               Département de Mathématiques-Informatique
% Année       : 2025-2026
% Description : Fichier principal de l'UE SIS589. Orchestre tous les chapitres,
%               les laboratoires, la bibliographie et les éléments de couverture.
%               Structuré en quatre parties thématiques + trois labs immersifs.
% Compilateur : lualatex + biber (recommandé)
%               → lualatex main && biber main && lualatex main && lualatex main
%               Alternative : pdflatex + biber (même séquence)
% ==============================================================================

\documentclass[12pt, a4paper, twoside, openright]{book}

% ==============================================================================
% CHARGEMENT DE LA FEUILLE DE STYLE
% ==============================================================================
\usepackage{style}

% ==============================================================================
% BIBLIOGRAPHIE
% ==============================================================================
\addbibresource{bibliographie.bib}

% ==============================================================================
% MÉTADONNÉES PDF
% ==============================================================================
\hypersetup{
    pdftitle   = {Supervision, Analyse Forensique et Gestion des Incidents},
    pdfauthor  = {MINKA MI NGUIDJOÏ Thierry Emmanuel},
    pdfsubject = {Support de Cours — SOC, SIEM, IDS, Forensics, IR, CERT},
    pdfkeywords= {SOC, SIEM, IDS, Forensics, ISO 27035, CSIRT, CERT,
                  Incident Response, Investigation Numérique, Cybersécurité,
                  Université de Yaoundé I, SIS589}
}

% ==============================================================================
% DÉBUT DU DOCUMENT
% ==============================================================================
\begin{document}

% ------------------------------------------------------------------------------
% 1. PAGE DE TITRE (1ÈRE DE COUVERTURE)
% ------------------------------------------------------------------------------
\begin{titlepage}
\thispagestyle{empty}
\pagecolor{CouleurPrimaire}

\begin{tikzpicture}[remember picture, overlay]
    % Bande accent en bas
    \fill[CouleurAccent]
        (current page.south west) rectangle
        ([yshift=2.2cm] current page.south east);
    % Ligne décorative en haut
    \fill[CouleurAccent]
        ([yshift=-1.5cm] current page.north west) rectangle
        ([yshift=-1.8cm] current page.north east);
    % Élément graphique — grille hex (évocation SOC/réseau)
    \fill[white, opacity=0.03]
        ([xshift=7.5cm, yshift=-5.5cm] current page.north west)
        -- ++(5.5,0) -- ++(0,-5.5) -- ++(-5.5,0) -- cycle;
    \fill[white, opacity=0.05]
        ([xshift=9cm, yshift=-7cm] current page.north west)
        -- ++(6.5,0) -- ++(0,-6.5) -- ++(-6.5,0) -- cycle;
    % Accent diagonal
    \fill[CouleurAccent, opacity=0.25]
        ([xshift=13cm, yshift=0cm] current page.south west)
        -- ++(3,0) -- ([xshift=0cm, yshift=3cm] current page.south east)
        -- ([xshift=0cm, yshift=0cm] current page.south east) -- cycle;
\end{tikzpicture}

\vspace*{1.5cm}

% Logo institutionnel (décommenter si disponible)
% \begin{center}
%     \includegraphics[height=2.5cm]{figures/logo_ensp.png}
% \end{center}

\begin{center}
    {\color{white!80!gray}\small\scshape
        Université de Garoua~·~Faculté des Sciences~·~%
        Département de Mathématiques-Informatique}\\[0.4ex]
    {\color{CouleurAccent}\rule{0.55\linewidth}{0.8pt}}
\end{center}

\vspace{1.5cm}

\begin{center}
    {\color{white!60!gray}\footnotesize\scshape support de cours}\\[1ex]
    {\color{white}\fontsize{30}{38}\selectfont\bfseries
        Supervision, Analyse Forensique}\\[0.6ex]
    {\color{white}\fontsize{30}{38}\selectfont\bfseries
        et Gestion des Incidents}\\[1.5ex]
    {\color{CouleurAccent}\Large\itshape
        SOC~·~SIEM~·~IDS~·~Forensics~·~Réponse aux Incidents}
\end{center}

\vspace{1.8cm}

\begin{center}
    {\color{white!70!gray}\small
        \ding{43}\quad Cas fil rouge : Entreprise LiZbiZ-SIS}\\[0.5ex]
    {\color{white!70!gray}\small
        \ding{44}\quad ISO 27035~·~ISO 27037~·~RFC 3227~·~NIST SP 800-86}\\[0.5ex]
    {\color{white!70!gray}\small
        \ding{228}\quad MITRE ATT\&CK~·~Kill Chain~·~Diamond Model~·~STIX/TAXII}\\[0.5ex]
    {\color{white!70!gray}\small
        \ding{228}\quad Loi Camerounaise 2010/012~·~Convention de Malabo}
\end{center}

\vfill

\begin{center}
    {\color{CouleurAccent}\rule{0.7\linewidth}{1.2pt}}\\[1.2ex]
    {\color{white}\large\bfseries MINKA MI NGUIDJOÏ Thierry Emmanuel}\\[0.4ex]
    {\color{white!80!gray}\small
        CISA~·~CISM~·~CGEIT~·~CRISC~·~CDPSE~·~ISO 27001~·~ISO 22301~·~ISO 9001}\\[0.3ex]
    {\color{white!70!gray}\footnotesize\itshape
        Chercheur au Laboratoire d'Ingénierie Mathématiques et des Systèmes d'Information}\\[0.3ex]
    {\color{white!70!gray}\footnotesize\itshape
        Expert Judiciaire~—~Analyste Forensique~—~Auditeur Senior~—~Cryptographe}\\[1.2ex]
    {\color{CouleurAccent}\rule{0.7\linewidth}{1.2pt}}
\end{center}

\vspace{0.6cm}

\begin{center}
    {\color{white!60!gray}\small Année académique 2025--2026}
\end{center}

\afterpage{\nopagecolor}
\end{titlepage}


% ------------------------------------------------------------------------------
% 2. PAGE DE DROITS / VERSO DE COUVERTURE
% ------------------------------------------------------------------------------
\clearpage
\thispagestyle{empty}
\vspace*{\fill}

\begin{center}
    \begin{minipage}{0.78\textwidth}
        \footnotesize\color{CouleurBordInfo}
        \begin{center}
            {\bfseries\color{CouleurPrimaire}
                Supervision, Analyse Forensique et Gestion des Incidents}\\[0.4ex]
            {\itshape SOC~·~SIEM~·~IDS~·~Forensics~·~Réponse aux Incidents}
        \end{center}
        \vspace{0.8ex}
        \hrule
        \vspace{0.8ex}
        \textbf{Auteur :} MINKA MI NGUIDJOÏ Thierry Emmanuel\\
        \textbf{Institution :} Université de Garoua — Faculté des Sciences — Département de Mathématiques-Informatique\\
        \textbf{Code UE :} SIS589 — Niveau~5\\
        \textbf{Année :} 2025--2026\\[0.5ex]
        Ce document est un support de cours universitaire rédigé à des fins
        pédagogiques. Toute reproduction partielle ou totale à des fins
        commerciales est interdite sans l'autorisation écrite de l'auteur.
        Les références à des produits, organisations ou normes sont citées
        à titre pédagogique uniquement.\\[0.5ex]
        Les certifications \textbf{CISA}, \textbf{CISM}, \textbf{CGEIT},
        \textbf{CRISC} et \textbf{CDPSE} sont des marques déposées de
        l'\textbf{ISACA}. \textbf{MITRE ATT\&CK} est une marque déposée
        de The MITRE Corporation. \textbf{EBIOS RM} est une méthode
        de l'\textbf{ANSSI}.\\[0.5ex]
        \hrule
        \vspace{0.5ex}
        {\centering\itshape
        «~Détecter sans comprendre, c'est voir sans regarder.
        Investiguer sans rigueur, c'est chercher sans trouver.~»}
    \end{minipage}
\end{center}

\vspace*{\fill}


% ------------------------------------------------------------------------------
% 3. DÉDICACE
% ------------------------------------------------------------------------------
\clearpage
\thispagestyle{empty}
\vspace*{3cm}
\begin{flushright}
    \begin{minipage}{0.62\textwidth}
        \itshape\large\color{CouleurPrimaire}
        À ceux qui, dans les salles obscures des SOC,\\
        veillent quand les autres dorment.\\[1.5ex]
        \normalsize\normalfont\color{CouleurAccent}
        Et à tous les enquêteurs numériques en devenir\\
        qui apprendront ici à lire les traces\\
        avant que le temps ne les efface.
    \end{minipage}
\end{flushright}


% ------------------------------------------------------------------------------
% 4. REMERCIEMENTS
% ------------------------------------------------------------------------------
\clearpage
\thispagestyle{empty}
\chapter*{Remerciements}
\addcontentsline{toc}{chapter}{Remerciements}

La conception de ce cours s'appuie sur l'expérience conjuguée du praticien
et du chercheur. L'auteur exprime sa profonde reconnaissance à tous les
relecteurs, anonymes ou non, qui ont contribué à enrichir et à affiner
chaque version de ce support.

Ce cours a été nourri par des années de missions opérationnelles en
investigation numérique, gestion d'incidents critiques et supervision
de systèmes d'information dans des organisations publiques et privées
au Cameroun et au-delà. La confrontation permanente entre théorie
et terrain donne à ces pages leur ancrage dans la réalité.

À nos étudiants du niveau~5, qui constitueront demain la première
ligne de défense de nos systèmes d'information nationaux.


% ------------------------------------------------------------------------------
% 5. PRÉAMBULE PÉDAGOGIQUE
% ------------------------------------------------------------------------------
\clearpage
\thispagestyle{empty}
\chapter*{Préambule}
\addcontentsline{toc}{chapter}{Préambule}

\begin{boxinfo}
\textbf{Prérequis recommandés :}
\begin{itemize}
    \item Réseaux informatiques : TCP/IP, protocoles, firewalls, segmentation VLAN.
    \item Systèmes d'exploitation : administration Linux et Windows (niveau opérationnel).
    \item Cours d'Audit des SI et/ou de Hacking Éthique fortement appréciés en préalable.
    \item Notions de droit numérique : RGPD, responsabilité civile et pénale, cybercriminalité.
\end{itemize}
\end{boxinfo}

\medskip

Ce cours aborde la sécurité de l'information sous l'angle de la \textbf{réactivité
opérationnelle} : détecter l'intrusion, en comprendre le déroulé, en préserver les
traces et en coordonner la réponse. Ces quatre actes — \textit{détecter, investiguer,
préserver, répondre} — constituent l'essence de la mission SIS589.

Le cas fil rouge \textbf{LiZbiZ-SIS} (extension de l'entreprise LiZbiZ introduite
dans les cours précédents) servira de terrain d'application systématique à chaque
chapitre. Les laboratoires en fin de parcours consolideront les compétences
sur environnements simulés (GNS3 / EVE-NG / SIFT Workstation).


% ------------------------------------------------------------------------------
% 6. TABLES
% ------------------------------------------------------------------------------
\clearpage
\tableofcontents

\clearpage
\listoffigures
\addcontentsline{toc}{chapter}{Liste des figures}

\listoftables
\addcontentsline{toc}{chapter}{Liste des tableaux}

\lstlistoflistings
\addcontentsline{toc}{chapter}{Liste des codes source}


% ==============================================================================
% CORPS DU COURS
% ==============================================================================
\mainmatter


% ------------------------------------------------------------------------------
% INTRODUCTION GÉNÉRALE
% ------------------------------------------------------------------------------
\input{chapitres/00_introduction}


% ============================================================
% PARTIE I — SUPERVISION ET DÉTECTION
% ============================================================
\part{Supervision et Détection}

% Chapitre 1 — Architecture SOC, métriques, gouvernance
\input{chapitres/01_supervision}

% Chapitre 2 — Logs : collecte, normalisation, intégrité, corrélation
\input{chapitres/02_logs}

% Chapitre 3 — IDS/IPS/NDR : signatures, anomalies, EDR, XDR
\input{chapitres/03_ids_ips}


% ============================================================
% PARTIE II — CADRE LÉGAL ET NORMATIF
% ============================================================
\part{Cadre Légal et Normatif}

% Chapitre 4 — ISO 27035, RFC 3227, NIST SP 800-86,
%              Loi 2010/012 du Cameroun, Convention de Malabo
\input{chapitres/04_cadre_legal}


% ============================================================
% PARTIE III — INVESTIGATION NUMÉRIQUE (FORENSICS)
% ============================================================
\part{Investigation Numérique}

% Chapitre 5 — Fondements : ordre de volatilité, chaîne de custody,
%              principes ACPO, admissibilité des preuves
\input{chapitres/05_forensics_fondements}

% Chapitre 6 — Forensique système : disque, système de fichiers,
%              artefacts Windows/Linux, timeline
\input{chapitres/06_forensics_systeme}

% Chapitre 7 — Forensique mémoire, réseau et cloud
\input{chapitres/07_forensics_memoire_reseau}


% ============================================================
% PARTIE IV — RÉPONSE AUX INCIDENTS ET AMÉLIORATION CONTINUE
% ============================================================
\part{Réponse aux Incidents et Amélioration Continue}

% Chapitre 8 — Cycle ISO 27035, CSIRT/CERT, Playbook, Runbook,
%              MITRE ATT&CK, qualification des incidents
\input{chapitres/08_incident_response}

% Chapitre 9 — Threat Intelligence : CTI, STIX/TAXII, IoC,
%              Kill Chain, Diamond Model
\input{chapitres/09_threat_intelligence}

% Chapitre 10 — RETEX, amélioration continue, métriques SOC
\input{chapitres/10_retex_amelioration}

% Chapitre 11 — Évaluation finale et grille de compétences
\input{chapitres/11_evaluation}


% ============================================================
% TRAVAUX PRATIQUES IMMERSIFS
% ============================================================
\part{Travaux Pratiques Immersifs}

% Lab 1 — Déploiement SIEM, corrélation d'événements, chasse aux alertes
\input{labs/lab1_siem_detection}

% Lab 2 — Acquisition d'image disque, analyse Autopsy / FTK Imager
\input{labs/lab2_forensics_disque}

% Lab 3 — Simulation Tabletop : scénario ransomware sur LiZbiZ-SIS
\input{labs/lab3_incident_simulation}


% ------------------------------------------------------------------------------
% CONCLUSION GÉNÉRALE
% ------------------------------------------------------------------------------
\input{chapitres/00_conclusion}


% ==============================================================================
% MATIÈRE POSTLIMINAIRE
% ==============================================================================
\backmatter

% --- Bibliographie ---
\clearpage
\printbibliography[
    heading=bibintoc,
    title={Bibliographie et Références}
]

% --- À Propos de l'Auteur ---
\clearpage
\input{annexes/about_author}

% --- Glossaire ---
\clearpage
% \printglossaries   % à activer après makeglossaries

% --- Index ---
\clearpage
\printindex
\addcontentsline{toc}{chapter}{Index}


% ==============================================================================
% ANNEXES
% ==============================================================================
\clearpage
\appendix

\input{annexes/tools}
\input{annexes/templates}
\input{annexes/glossary}


% ==============================================================================
% 4ÈME DE COUVERTURE
% ==============================================================================
\clearpage
\thispagestyle{empty}
\pagecolor{CouleurPrimaire}

\begin{tikzpicture}[remember picture, overlay]
    \fill[CouleurAccent]
        (current page.south west) rectangle
        ([yshift=1.8cm] current page.south east);
    \fill[CouleurAccent]
        ([yshift=-1.4cm] current page.north west) rectangle
        ([yshift=-1.7cm] current page.north east);
\end{tikzpicture}

\vspace*{2.5cm}

\begin{center}
    {\color{CouleurAccent}\rule{0.6\linewidth}{1pt}}\\[1.5ex]
    {\color{white}\Large\bfseries À propos de ce cours}\\[1.5ex]
    {\color{CouleurAccent}\rule{0.6\linewidth}{1pt}}
\end{center}

\vspace{1cm}

\begin{center}
\begin{minipage}{0.75\textwidth}
\color{white!85!gray}\setstretch{1.4}
Ce cours constitue le maillon opérationnel de la chaîne de compétences en
sécurité des systèmes d'information. Là où l'audit identifie et qualifie les
risques, SIS589 forme à \textbf{détecter}, \textbf{investiguer} et \textbf{maîtriser}
les incidents en temps réel.

\medskip

Articulé autour du triptyque \textit{SOC~/ Forensics~/ Incident Response},
il couvre l'intégralité du cycle opérationnel~: supervision continue,
analyse de logs, investigation numérique légale, réponse structurée aux
incidents, renseignement sur la menace et retour d'expérience.

\medskip

Ancré dans les standards internationaux (ISO 27035, RFC 3227,
NIST SP 800-86, MITRE ATT\&CK) et le cadre juridique camerounais
(Loi 2010/012, Convention de Malabo), ce support s'adresse aux
ingénieurs et spécialistes en cybersécurité de niveau~5 appelés à
opérer dans des environnements réels et sous contrainte juridique.
\end{minipage}
\end{center}

\vfill

\begin{center}
    {\color{CouleurAccent}\rule{0.7\linewidth}{1.2pt}}\\[1.2ex]
    {\color{white}\large\bfseries MINKA MI NGUIDJOÏ Thierry Emmanuel}\\[0.4ex]
    {\color{white!80!gray}\small
        CISA~·~CISM~·~CGEIT~·~CRISC~·~CDPSE~·~ISO 27001~·~ISO 22301~·~ISO 9001}\\[0.3ex]
    {\color{white!70!gray}\footnotesize\itshape
        Chercheur au laboratoire d'Ingénierie Mathématiques et des Systèmes d'Information}\\[0.3ex]
    {\color{white!70!gray}\footnotesize\itshape
        Expert Judiciaire~—~Analyste Forensique~—~Auditeur Senior~—~Cryptographe}\\[1.2ex]
    {\color{CouleurAccent}\rule{0.7\linewidth}{1.2pt}}
\end{center}

\vspace{1cm}

\afterpage{\nopagecolor}

\end{document}
% ==============================================================================
% Fin de main.tex — UE SIS589
% ==============================================================================
